\cotefile{138}
\scanpage{32}

\deleted{$\widetilde E_1$\unsure{fixe, $\alpha$|$\{\alpha,\alpha\}$}}

Soit

\[
\deleted{\widetilde C_1=\bigl\{(\zeta_1,\zeta_2,\zeta_\infty)\in\mu_3^3}
\]

\[
\begin{aligned}
\widetilde C_1&=\operatorname{Isom}(\{\underbrace{1,2,\infty}_{I}\},\mu_3)
=\bigl\{(\zeta_1,\zeta_2,\zeta_3)\in\mu_3^3\mid\\
&\hspace{17em}\zeta_1\ne\zeta_2\ne\zeta_3\ne\zeta_1\bigr\}\\[-1ex]
&\subset\deleted{\unsure{\overline{\mathbb Q}^{\,3}}}\ \text{cône projectant de }\overline C_1\subset\overline{\mathbb Q}^{\,3}
\end{aligned}\tag*{(24)}
\]

$\widetilde C_1$ a $3!=6$ éléments, et $\mu_3$ y opère diagonalement \deleted{sans point fixe} \inserted{librement (multiplication scalaire)}, le quotient $\widetilde C_1/\mu_3$ a donc deux éléments. L'application

\[
\widetilde C_1/\mu_3\longrightarrow\overline C_1
\]

\margin{\uline{NB} $\widetilde C_1/\mu_3\subset\Pi_3$ est formé des \deleted{\unsure{liens|lignes}}}

induit une bijection

\[
\widetilde C_1/\mu_3\simeq\overline C_1^{\unsure{\mathfrak S_3|\mathscr S_3}^{+}}
=\{\unsure{\sigma|\alpha|\varpi},\bar\sigma\}
\tag*{(25)}
\]

mettant en évidence l'opération de $\mathcal E$ sur $(\sigma,\bar\sigma)$, via ses opérations sur $\widetilde C_1\subset\mu_3^3=\mu_3^I$ (\uline{NB} $\mathcal E$ opère via $\mathfrak S_3\times\{\pm1\}$, avec $\Gamma\mathrel{\overset{\chi_3}{\unsure{-|\longrightarrow}}}\{\pm1\}$)

Soit alors $\widetilde C_n$ l'image inverse de $\widetilde C_1$ par l'application allant du cône projectant de $\overline C_n$, au cône projectant de $\overline C_1$, donc \deleted{c'est l'ens de $(x,y,z)$ avec \unsure{$z$}}

\[
\left\{\begin{aligned}
\overline C_n&=\{(x,y,z)\in\overline{\mathbb Q}^{\,3}\mid
\deleted{x^n+y^n+z^n=1,\ x}\inserted{(x^n,y^n,z^n)\in\widetilde C_1}\}\\
&=\{(x;y,z)\in\overline{\mathbb Q}^{\,3}\mid
x^{3n}=y^{3n}=z^{3n}=1,\ x^{\unsure{3n|n}}\ne y^n\ne z^n\ne x^n\}
\end{aligned}\right.\tag*{(26)}
\]

Noter que $\mu_{3n}$ opère diagonalement sur $\overline C_n$,

% archive sheet number: 31
